% Options for packages loaded elsewhere
\PassOptionsToPackage{unicode}{hyperref}
\PassOptionsToPackage{hyphens}{url}
%
\documentclass[
]{scrbook}
\usepackage{amsmath,amssymb}
\usepackage{iftex}
\ifPDFTeX
  \usepackage[T1]{fontenc}
  \usepackage[utf8]{inputenc}
  \usepackage{textcomp} % provide euro and other symbols
\else % if luatex or xetex
  \usepackage{unicode-math} % this also loads fontspec
  \defaultfontfeatures{Scale=MatchLowercase}
  \defaultfontfeatures[\rmfamily]{Ligatures=TeX,Scale=1}
\fi
\usepackage{lmodern}
\ifPDFTeX\else
  % xetex/luatex font selection
\fi
% Use upquote if available, for straight quotes in verbatim environments
\IfFileExists{upquote.sty}{\usepackage{upquote}}{}
\IfFileExists{microtype.sty}{% use microtype if available
  \usepackage[]{microtype}
  \UseMicrotypeSet[protrusion]{basicmath} % disable protrusion for tt fonts
}{}
\makeatletter
\@ifundefined{KOMAClassName}{% if non-KOMA class
  \IfFileExists{parskip.sty}{%
    \usepackage{parskip}
  }{% else
    \setlength{\parindent}{0pt}
    \setlength{\parskip}{6pt plus 2pt minus 1pt}}
}{% if KOMA class
  \KOMAoptions{parskip=half}}
\makeatother
\usepackage{xcolor}
\usepackage{longtable,booktabs,array}
\usepackage{calc} % for calculating minipage widths
% Correct order of tables after \paragraph or \subparagraph
\usepackage{etoolbox}
\makeatletter
\patchcmd\longtable{\par}{\if@noskipsec\mbox{}\fi\par}{}{}
\makeatother
% Allow footnotes in longtable head/foot
\IfFileExists{footnotehyper.sty}{\usepackage{footnotehyper}}{\usepackage{footnote}}
\makesavenoteenv{longtable}
\usepackage{graphicx}
\makeatletter
\def\maxwidth{\ifdim\Gin@nat@width>\linewidth\linewidth\else\Gin@nat@width\fi}
\def\maxheight{\ifdim\Gin@nat@height>\textheight\textheight\else\Gin@nat@height\fi}
\makeatother
% Scale images if necessary, so that they will not overflow the page
% margins by default, and it is still possible to overwrite the defaults
% using explicit options in \includegraphics[width, height, ...]{}
\setkeys{Gin}{width=\maxwidth,height=\maxheight,keepaspectratio}
% Set default figure placement to htbp
\makeatletter
\def\fps@figure{htbp}
\makeatother
\setlength{\emergencystretch}{3em} % prevent overfull lines
\providecommand{\tightlist}{%
  \setlength{\itemsep}{0pt}\setlength{\parskip}{0pt}}
\setcounter{secnumdepth}{5}
% definitions for citeproc citations
\NewDocumentCommand\citeproctext{}{}
\NewDocumentCommand\citeproc{mm}{%
  \begingroup\def\citeproctext{#2}\cite{#1}\endgroup}
\makeatletter
 % allow citations to break across lines
 \let\@cite@ofmt\@firstofone
 % avoid brackets around text for \cite:
 \def\@biblabel#1{}
 \def\@cite#1#2{{#1\if@tempswa , #2\fi}}
\makeatother
\newlength{\cslhangindent}
\setlength{\cslhangindent}{1.5em}
\newlength{\csllabelwidth}
\setlength{\csllabelwidth}{3em}
\newenvironment{CSLReferences}[2] % #1 hanging-indent, #2 entry-spacing
 {\begin{list}{}{%
  \setlength{\itemindent}{0pt}
  \setlength{\leftmargin}{0pt}
  \setlength{\parsep}{0pt}
  % turn on hanging indent if param 1 is 1
  \ifodd #1
   \setlength{\leftmargin}{\cslhangindent}
   \setlength{\itemindent}{-1\cslhangindent}
  \fi
  % set entry spacing
  \setlength{\itemsep}{#2\baselineskip}}}
 {\end{list}}
\usepackage{calc}
\newcommand{\CSLBlock}[1]{\hfill\break\parbox[t]{\linewidth}{\strut\ignorespaces#1\strut}}
\newcommand{\CSLLeftMargin}[1]{\parbox[t]{\csllabelwidth}{\strut#1\strut}}
\newcommand{\CSLRightInline}[1]{\parbox[t]{\linewidth - \csllabelwidth}{\strut#1\strut}}
\newcommand{\CSLIndent}[1]{\hspace{\cslhangindent}#1}
\usepackage{booktabs} % for better tables

% fonts
% libertine font: https://www.tug.org/FontCatalogue/linuxlibertine/
\usepackage[oldstyle]{libertine}
\usepackage{libertinust1math}
% \usepackage[T1]{fontenc}
\usepackage[extralight]{inter} 

\usepackage[left=3cm, right=2.5cm, top=2.5cm, bottom=4cm]{geometry} % page margings

\usepackage{tabto} % for aligning text at a certain position, used for title page

\usepackage[ngerman, english]{babel}		% main language is the last in the list, more infor is here: https://en.wikibooks.org/wiki/LaTeX/Internationalization#Babel

\usepackage[onehalfspacing]{setspace}	% clean 1.5 leading / spacing
\RedeclareSectionCommand[beforeskip=0.5em,afterskip=2em]{chapter} % define space before and after heading

\usepackage{pdfpages} % to include pdfs
\ifLuaTeX
  \usepackage{selnolig}  % disable illegal ligatures
\fi
\usepackage{bookmark}
\IfFileExists{xurl.sty}{\usepackage{xurl}}{} % add URL line breaks if available
\urlstyle{same}
\hypersetup{
  hidelinks,
  pdfcreator={LaTeX via pandoc}}

\author{}
\date{\vspace{-2.5em}}

\begin{document}

\frontmatter

\begin{titlepage}                   
    \begin{center}
        \includegraphics[height=3.2cm]{logo.pdf}\\[40mm]
        
        \huge {\linespread{1.5} Rethinking Variation in Social Cognition: \\ Gaze Following across Individuals, Ages, and Communities}\\[30mm]
        
                
        \normalsize Von der Fakultät Nachhaltigkeit \\ der Leuphana Universität Lüneburg zur Erlangung des Grades\\
        Doktorin der Psychologie\\-- Dr. rer. nat. --\\[10mm]
        
        genehmigte Dissertation von\\[10mm]
        Julia Christin Prein\\
        geboren am 26.02.1995 in Berlin-Kreuzberg, Deutschland
        
    \vspace*{\fill} 
    \end{center}
    
    \newpage
    \thispagestyle{empty}
    \begin{flushleft}
      \begin{normalsize}
      
      \vspace*{\fill} 
            Eingereicht am: \tabto*{30mm} 1. Oktober 2024 \\[10mm]
            
            Erstbetreuer: \tabto*{30mm} Prof.\,Dr.\, Manuel Bohn, \textit{Leuphana Universität Lüneburg}\\
                Zweitbetreuer: \tabto*{30mm} Prof.\,Dr.\, Sebastian Wallot, \textit{Leuphana Universität Lüneburg}\\[10mm]
                
                Erstgutachter: \tabto*{30mm} Prof.\,Dr.\, Manuel Bohn, \textit{Leuphana Universität Lüneburg}\\
                Zweitgutachter: \tabto*{30mm} Prof.\,Dr.\, Sebastian Wallot, \textit{Leuphana Universität Lüneburg}\\
                Drittgutachter: \tabto*{30mm} Prof.\,Dr.\, Daniel Haun, \textit{Universität Leipzig \& Max-Planck-Institut für\\
                \tabto*{30mm} Evolutionäre Anthropologie}\\[10mm]
        \end{normalsize}
    \end{flushleft}
    
\end{titlepage}

\tableofcontents

\chapter{Acknowledgments}\label{acknowledgments}

This dissertation rests on the shoulders of many inspiring, thought-provoking, motivating, and caring individuals. The last four years have been many different things --- but definitely neither boring nor stagnant. I want to thank everyone who has accompanied me on this journey.

My motivation to work on this dissertation is greatly due to my wonderful supervisors, Manuel and Daniel. Manuel, I cannot thank you enough for supporting and believing in me. You have taught me to critically challenge the status quo in research and to focus my energy on improving it --- pragmatically, step by step. Even though I do know that you believe there \emph{are} stupid questions, I felt very comfortable to ask you all along the way and I have benefited a lot from your patience with me. Thanks for roughly 200 meetings that have always left me more motivated and focused than I was before. I cannot imagine how science would be without you! Daniel, I am so happy and proud to have been a member of the CCP team. Thank you for offering me so many unique opportunities and unforgettable memories to travel to conferences and Uganda. I am deeply impressed by how quickly you get a grasp of a new project and know immediately which questions to ask. Thank you for being such an approachable director and having an open door to hear about my personal challenges. To my external supervisor, Robert: your teaching in Research Methods has turned my hate for statistics into a deep passion. Thanks for looking into open PhD positions with and for me, and accompanying me along the way. Without you, I would have never even considered doing a PhD.

Thanks to the new people that I met at the Leuphana. Sebastian, thank you for making it possible for me to pursue this doctoral title. Patricia, Franziska, Nele, and (again) Manuel, thanks for the welcoming and kind atmosphere in the LüneLütten lab. I am excited about what is to come next!

I am grateful to our wonderful cohort of PhDs at the MPI. You are a remarkable group of people, and I am very happy that we could share our daily struggles with each other. I will miss seeing you every day! Noemi, I am so glad that we got to share our office for a (way too short) period of time and successfully solved the Christmas murder mystery with Marie. Wilson, thanks for enduring our laughter through the office walls, for countless lunch breaks together, and for sharing random stories about your fieldwork (and Iceland)! Without my work wife Marie, my time during this dissertation would not have been even half as great as it has been. I could (and have and will) cry thinking about not sharing my daily life with you anymore --- even though I completely trust in our relationship to live until we die. You have made me laugh so much, even if all I wanted to do was cry. Whatever it is that I am up to, I know you will always have my back and help me get back on my feet if I fail. I love your honesty, creativity, immense courage, and charming awkwardness, and feel so honored that you are my friend. Thank you for being you!

My dear, dear family, I am so deeply grateful to know you by my side. Whenever I talk about individual differences, I think about you, Caro, and how we two sisters have the same parents and yet, are so different (in the very best way!!). I believe it is rare that siblings do not feel the need to compete with each other but can just love each other as they are. Thank you. To my mom, who told me right in the beginning that this journey would be a marathon and not a sprint. In some of my darker moments, you told me that even if it felt like I was about to fall, I could never --- because there are so many caring people who would catch me beforehand. Thank you for always having an open home and lots of support (in the form of love, food, vacations, phone calls, and advice) for me. To my dad, thanks for showing up to Leipzig so often and establishing a little concert routine together! I believe I have learned a lot of my problem-solving skills and patience from you. Even if you do not remember saying ``Mach Dinge richtig'' in Caro's and my childhood, I think my dissertation has greatly benefited from this sentence. Thank you, Louis, for helping with my move to Leipzig and spontaneously catching up over coffee in Leipzig or Berlin. To my grandpa, who told me he never worried that I would find my own way. And to my grandma, for all your curiosity, play, and love for children. To the whole Gerken-Prein clan (including our furry best friends Resi and Milou, reminding me to cuddle and catch some fresh air), I could not wish for a better family. Finally, we are all in the North again, reunited.

Thank you to all my amazing friends, who constantly remind me to live my life outside the office. To O74, OLAV, and Shubhangi, Schlotti, and Maló, for the cozy home we built together and all the little dance breaks, kitchen karaoke, reading clubs, and your patience with my mood swings and stealing lunch boxes. With you, even COVID isolation was fun. Thanks to Laura, for drying my tears, always having an open ear and heart, and helping with every single move. To Marvin, for all the long hours of deep talk and dancing. Thanks to Arne, for nerding about statistics and providing way too much green mouth refreshener. To the Osna Coxis, especially Merle, Leah, and Alissa, for all the New Year's Eves, regular video calls, big hugs and words of comfort. To Niki, for constantly reaching out, joking around, and asking me how I'm doing. To Ani, Nica, and Schlotti for a lifelong friendship, wholehearted laughter, and a traveling diary. I am so grateful to have you all in my life.

And finally, Steven. You have taken over so many different roles in my life over the last four years. In the beginning, I was simply happy to get some technological support and coding enthusiasm from you. Thank you for pushing me beyond my levels of comfort and believing in my programming skills (\ldots{} and for being my IT partner ;)). Quickly, however, coding became only one of the thousand topics that we would regularly discuss. Every morning, I looked forward to you making creepy, funny faces on our office door and having long, controversial, and nerdy conversations with you. I deeply miss you in my daily work life. And yet, I do appreciate all the changes that our relationship has gone through in recent years --- through all the highs and lows. I am so sincerely grateful and glad to share my life with you. It is so easy being myself around you. Thank you for all the commitment, stability, honesty, relaxation, music, joy, and silly jokes that you bring into my life. Your love and support means the world to me. I am excited about our future in Hamburg!

\chapter{Copyright Notice}\label{copyright}

The research articles included in this cumulative dissertation have been or will be published in international peer-reviewed journals. Copyright of the text and illustrations lies with the author or authors of the respective chapter(s). The publishers own the exclusive rights to publish or use the text and illustrations for their own purposes. Reprinting any part of this dissertation thesis requires the permission of the copyright holder(s).
The individual contributions of this cumulative thesis have been or will be published (in chronological order) as follows:

Prein, J. C., Kalinke, S., Haun, D. B. M.*, \& Bohn, M.* (2023). TANGO: A reliable, open-source, browser-based task to assess individual differences in gaze understanding in 3 to 5-year-old children and adults. \emph{Behavior Research Methods, 56}(3), 2469--2485. \url{https://doi.org/10.3758/s13428-023-02159-5}

Prein, J. C., Maurits, L., Werwach, A., Haun, D. B. M.,* \& Bohn, M.* (2024). \emph{Variation in gaze understanding across the life span: A process-level perspective.} PsyArXiv. \url{https://doi.org/10.31234/osf.io/dy73a}

Bohn, M.*, Prein, J. C.*, Ayikoru, A., Bednarski, F. M., Dzabatou, A., Frank, M. C., Henderson, A. M. E., Isabella, J., Kalbitz, J., Kanngiesser, P., Keşşafoğlu, D., Koymen, B., Manrique-Hernandez, M., Magazi, S., Mújica-Manrique, L., Ohlendorf, J., Olaoba, D., Pieters, W., Pope-Caldwell, S., \ldots{} Haun, D. (2024). \emph{A universal of human social cognition: Children from 17 communities process gaze in similar ways} {[}Manuscript submitted for publication{]}. PsyArXiv. \url{https://doi.org/10.31234/osf.io/z3ahv}

Prein, J.C., Bednarski, F. M., Dzabatou, A., Frank, M. C., Henderson, A. M. E., Kalbitz, J., Kanngiesser, P., Keşşafoğlu, D., Köymen, B., Manrique-Hernandez, M. V., Magazi, S., Mújica-Manrique, L., Ohlendorf, J., Olaoba, D., Pieters, W. R., Pope-Caldwell, S., Sen, U., Slocombe, K., Sparks, R. Z., Stengelin, R., Sunderarajan, J., Sutherland, K., Tusiime, F., Vieira, W., Zhang, Z., Zong, Y., Haun, D. B. M.*, Bohn, M.*. (2024). \emph{Measuring variation in gaze following across communities, ages, and individuals -- a showcase of the TANGO--CC} {[}Manuscript submitted for publication{]}. PsyArxiv. \url{https://doi.org/10.31234/osf.io/fcq2g}

~

Further articles that were written in the context of this dissertation can be found in the Appendix and have been or will be published as follows (in chronological order):

Schuwerk, T., Kampis, D., Baillargeon, R., Biro, S., Bohn, M., Byers-Heinlein, K., Dörrenberg, S., Fisher, C., Franchin, L., Fulcher, T., Garbisch, I., Geraci, A., Grosse Wiesmann, C., Hamlin, K., Haun, D. B. M., Hepach, R., Hunnius, S., Hyde, D. C., Karman, P., \ldots, Prein, J., \ldots{} Rakoczy, H. (2021). \emph{Action anticipation based on an agent's epistemic state in toddlers and adults. Child Development} {[}In-Principle Acceptance of Registered Report Stage 1: Study Design{]}. PsyArXiv. \url{https://doi.org/10.31234/osf.io/x4jbm}

Bohn, M., Prein, J. C., Engicht, J., Haun, D., Gagarina, N., \& Koch, T. (2023). \emph{PREVIC: An adaptive parent report measure of expressive vocabulary in children between 3 and 8 years of age} {[}Manuscript submitted for publication{]}. PsyArXiv. \url{https://doi.org/10.31234/osf.io/hvncp}

Bohn, M.*, Prein, J.*, Koch, T., Bee, R. M., Delikaya, B., Haun, D., \& Gagarina, N. (2024). oREV: An item response theory-based open receptive vocabulary task for 3- to 8-year-old children. \emph{Behavior Research Methods, 56}(3), 2595--2605. \url{https://doi.org/10.3758/s13428-023-02169-3}

Steffan, A., Zimmer, L., Arias-Trejo, N., Bohn, M., Dal Ben, R., Flores-Coronado, M. A., Franchin, L., Garbisch, I., Grosse Wiesmann, C., Hamlin, J. K., Havron, N., Hay, J. F., Hermansen, T. K., Jakobsen, K. V., Kalinke, S., Ko, E.-S., Kulke, L., Mayor, J., Meristo, M., \ldots, Prein, J., \ldots, Schuwerk, T. (2024). Validation of an open source, remote web-based eye-tracking method (WebGazer) for research in early childhood. \emph{Infancy, 29}(1), 31--55. \url{https://doi.org/10.1111/infa.12564}

~

* denotes shared first or last authorship.

~
~
~
~

\begingroup
\renewcommand{\cleardoublepage}{}
\renewcommand{\clearpage}{}

\chapter{Disclaimer}\label{disclaimer}

In this dissertation, the plural pronoun ``we'' highlights the collaborative efforts within the presented scientific projects. The contributions of all co-authors of the included manuscripts are transparently and accurately attributed using the Contributor Roles Taxonomy (CRediT; \url{https://credit.niso.org}) and can be found for each respective manuscript in {[}Appendix A{]}. In situations where the primary responsible person is not explicitly defined, ``we'' refers to the efforts of this dissertation's author alone, as the plural pronoun ``we'' reflects standard practice in psychological scientific writings. This use does not diminish the author's significance in conceptualizing, executing, documenting, and interpreting the presented research. The singular pronoun ``I'' explicitly denotes personal opinions, reflections, and decisions regarding this dissertation's theoretical framework paper and its structure.
\endgroup

\chapter{Abstract}\label{abstract}

Prein et al. (\citeproc{ref-prein2023tango}{2023})

Lorem ipsum dolor sit amet, consectetur adipiscing elit, sed do eiusmod tempor incididunt ut labore et dolore magna aliqua. Luctus venenatis lectus magna fringilla urna. Laoreet non curabitur gravida arcu ac tortor. Faucibus vitae aliquet nec ullamcorper sit amet risus nullam eget. Cras semper auctor neque vitae tempus quam pellentesque. Imperdiet massa tincidunt nunc pulvinar sapien et ligula ullamcorper. Et tortor at risus viverra adipiscing at. Congue nisi vitae suscipit tellus mauris. Habitant morbi tristique senectus et netus et malesuada fames ac. Eget mauris pharetra et ultrices neque. Aenean et tortor at risus viverra. Tempor orci dapibus ultrices in iaculis nunc sed augue. Euismod lacinia at quis risus sed vulputate odio ut enim. Id eu nisl nunc mi ipsum faucibus. Est lorem ipsum dolor sit amet. Eget velit aliquet sagittis id consectetur purus. Faucibus ornare suspendisse sed nisi. Pellentesque diam volutpat commodo sed egestas egestas fringilla.

\chapter{Zusammenfassung}\label{zusammenfassung}

Lorem ipsum dolor sit amet, consectetur adipiscing elit, sed do eiusmod tempor incididunt ut labore et dolore magna aliqua. Luctus venenatis lectus magna fringilla urna. Laoreet non curabitur gravida arcu ac tortor. Faucibus vitae aliquet nec ullamcorper sit amet risus nullam eget. Cras semper auctor neque vitae tempus quam pellentesque. Imperdiet massa tincidunt nunc pulvinar sapien et ligula ullamcorper. Et tortor at risus viverra adipiscing at. Congue nisi vitae suscipit tellus mauris. Habitant morbi tristique senectus et netus et malesuada fames ac. Eget mauris pharetra et ultrices neque. Aenean et tortor at risus viverra. Tempor orci dapibus ultrices in iaculis nunc sed augue. Euismod lacinia at quis risus sed vulputate odio ut enim. Id eu nisl nunc mi ipsum faucibus. Est lorem ipsum dolor sit amet. Eget velit aliquet sagittis id consectetur purus. Faucibus ornare suspendisse sed nisi. Pellentesque diam volutpat commodo sed egestas egestas fringilla.

\chapter{List of Abbreviations}\label{acronyms_HEADER_LOA}

\begin{description}
\tightlist
\item[\phantomsection\label{acronyms_ToM}{ToM}]
Theory of Mind
\end{description}

\mainmatter

\chapter{General Introduction}\label{introduction}

Test \hyperref[acronyms_ToM]{Theory of Mind (ToM)}

\section{Structure of this Dissertation}\label{structure-of-this-dissertation}

This dissertation is embedded in an overarching, still ongoing project led by Prof.~Dr.~Manuel Bohn (Leuphana University Lüneburg) and Prof.~Dr.~Daniel Haun (Max Planck Institute for Evolutionary Anthropology Leipzig) that will create a new task battery for social cognition. This task battery aims to capture individual differences in children from diverse communities. The present thesis describes the steps that have already been undertaken in this endeavor: during my time as a doctoral candidate, we have designed, implemented, and validated the first task of the battery. We have decided to focus the first task on one of the most fundamental social-cognitive abilities: gaze following.

The present dissertation seeks to approach the development of our new gaze following paradigm from two sides: from a contentual and a methodological perspective. The first part will focus on the psychological construct(s) underlying gaze following: I will describe the encompassing term of social cognition and its milestones during child development. I will continue to elaborate on the construct of gaze following, its current state of research and remaining open questions. Within both of these sections, I will consider existing studies that focus on individual differences and subsequently, studies that were conducted cross-culturally.

The second part will focus on the methodological considerations for constructing a new task battery. I will examine existing social cognition measures, identify common challenges and shortcomings, and establish the need for an individual differences- and cross-cultural perspective.

The remaining chapters will focus on our approach to develop a new gaze following measure that tackles these recognized methodological challenges. I will summarize our study findings, discuss their limitations and provide an outlook for future studies in social cognition, focusing on individual differences and cross-cultural research.

\section{Social Cognition}\label{social-cognition}

What makes us uniquely human? Roughly 90 percent of our DNA matches that of cats (\citeproc{ref-pontius2007initial}{Pontius et al., 2007}), and even close to 99 percent matches our closest living relatives, the chimpanzees (\citeproc{ref-waterson2005initial}{Waterson et al., 2005}). Consequently, we share several key features with each other. For example, our eyes, ears, and nose are designed to sensorily perceive the surroundings, and our neocortex processes this information and performs higher-order brain functions like spatial reasoning and motor commands. These features characterize all mammals. And yet, the majority of the earth's surface gets populated and exploited by mainly one mammal: the homo sapiens.

That it is humans who dominate the globe to such a degree (and even domesticate other mammals) is estonishing given the fact that there are many examples in which the animal kingdom far exceeds humans in their abilities. Cheetahs run with up to 120 km/h in speed, and Peregrine falcons reach over 380 km/h when diving for prey. Bats and dolphins use echolocation to navigate in complete darkness. Ocean Sunfish have extraordinary reproductive abilities and can lay up to 300 million eggs at once. The Greenland Shark can live as long as 400 years. So many other species are faster, produce more offspring, and live longer than humans, so which abilities brought humans to their ecological success?

One aspect that is often highlighted is our ability to engage in social interactions and communicate with others (\citeproc{ref-tomasello2019becoming}{Tomasello, 2019}). Humans have been described as ``ultra-social primates'' (\citeproc{ref-grossmann2022becoming}{Grossmann, 2022, p. 1}): we form social bonds (\citeproc{ref-bowlby1958nature}{Bowlby, 1958}), live (mostly) cooperatively (\citeproc{ref-tomasello2020adaptivea}{Tomasello, 2020}), exchange information and learn through observing others (\citeproc{ref-csibra2009natural}{Csibra \& Gergely, 2009}). By relying on these strategies, knowledge and skills could be passed on from generation to generation --- which could not have possibly been acquired by one single individual (\citeproc{ref-henrich2015secreta}{Henrich, 2015}).

A seminal paper by Herrmann et al. (\citeproc{ref-herrmann2007humansa}{2007}) tested the hypothesis that humans. They administered a cognitive battery with 16 different tasks to chimpanzees, orangutans, and 2.5-year-old human infants. Indeed, the researchers found that all three groups showed similar cognitive abilities within the physical world (space, quantities, causality), while human infants outperformed the chimpanzees and orangutans when it came to the social world (social learning, communication, Theory of Mind). This study clearly underlines humans' remarkable abilities in the social realm, which might explain why so many researchers focus on this domain: A quick search on the APA PsycNet (on June 25, 2024, at 17:00 CET) revealed 28,594 results for the search term ``social cognition'', while only 3,095 results appeared for ``spatial cognition'' and 95 appeared for ``physical cognition''.

\subsection{Terminology}\label{terminology}

Before going deeper into the developmental literature on social cognition, let us clarify the term: What is social cognition? Interestingly, there is not one agreed-upon definition accepted by most psychologists. Rather, social cognition is and cynical tongues might say that there are as many definitions of social cognition as there are research articles on it. As Ostrom put it in the Foreword of the Handbook of Social Cognition: ``I regard single-sentence definitions of social cognition to be slightly offensive {[}\ldots{]}. It would be more accurate to say that my preferred definition is the entire Handbook of Social Cognition, both this and the first edition combined.'' (\citeproc{ref-ostrom1994foreword}{Ostrom, 1994, p. vii}). Indeed, researchers have argued that the non-specific and heterogeneous vocabulary hinders empirical and theoretical advancement: sometimes, a single concept is described by several different labels (convergence of meaning), while other concepts share the same label but refer to different constructs (divergence of meaning) (\citeproc{ref-quesque2024defining}{Quesque et al., 2024}; \citeproc{ref-quesque2020theoryofmind}{Quesque \& Rossetti, 2020}).

To get a clearer grasp on what developmental psychologists understand under the term social cognition, we conducted a short expert survey at the beginning of my Ph.D.~(in Autumn/Winter 2020). We distributed an online questionnaire via personal contacts and emailing lists (e.g.~cogdevsoc mailing list) to researchers broadly working in cognitive, comparative, cross-cultural and/or developmental psychology. Out of the 100 experts that completed the survey, more than half held a professorship, and the majority focused on basic research with children. First, we asked experts to define social cognition as a psychological construct. Exactly one third of the sample (33\%, n = 33) chose a definition by Glynn \& Watkiss (\citeproc{ref-glynn2016social}{2016}): ``Social cognition is the process by which actors, at individual or collective levels, decode and encode their social world, using mental models, knowledge structures and cultural understandings to process information, extract meaning and determine appropriate action.'' (p.~1). The next two most frequently chosen definitions focused similarly on the information processing aspects in environments with other agents (``{[}\ldots{]} encompasses all the information-processing mechanisms that underlie how people capture, process, store, and apply information about others to navigate social situations'' (\citeproc{ref-decety2020preface}{Decety, 2020, p. ix}); ``{[}\ldots{]} is concerned with the study of the thought processes, both implicit and explicit, through which humans attain understanding of self, others, and their environment'' (\citeproc{ref-moskowitz2013social}{Moskowitz, 2013, p. 1})). In the next step, we asked experts to name which dimensions they mentioned most often when discussing or writing about social cognition. The dimensions that were chosen by more than half of the experts were beliefs, knowledge, perspective-taking and intentions. Finally, we compiled a list of 21 social-cognitive abilities and asked experts to rate the likelihood that children of the same age would vary to perform these abilities. The experts assumed that children varied mostly in whether they manipulate emotions, deceive others, simulate others' reasoning processes, take another's perspective, and understand the subjectivity of knowledge states. Abilities that were assumed to vary the least between children of the same age were recognizing agents and following gaze. It appeared that experts assumed variation in many social-cognitive abilities and judged the more room for individual differences, the more complex a given social-cognitive ability. Please note that more information on the expert survey can be found in {[}Appendix C{]}.

\begin{itemize}
\tightlist
\item
  distinction theory of mind vs social cognition needed?
\item
  ToM Abkürzung einführen
\end{itemize}

A related term to social cognition is \hyperref[acronyms_ToM]{ToM}. While these terms are sometimes used synonymously, social cognition functions as the umbrella term that encompasses \hyperref[acronyms_ToM]{ToM}. \hyperref[acronyms_ToM]{ToM} itself describes a system organized around beliefs, desires, and actions that we use in order to make sense of other agents. We try to understand actions by attributing mental states to ourselves and others --- summarized by Wellman as follows: ``in our everyday thinking we construe people as engaging in acts they \emph{think} will get them what they \emph{want}'' (\citeproc{ref-wellman2013universal}{Wellman, 2013, p. 69}).

\subsection{Developmental Milestones}\label{developmental-milestones}

Importantly, human infants do not come into the world already possessing all these social-cognitive abilities (\citeproc{ref-tomasello2020adaptivea}{Tomasello, 2020}). Much of the developmental psychology research today focuses on how children acquire the social-cognitive abilities needed to become full members of our society.

\begin{itemize}
\tightlist
\item
  developmental milestones in social cognition anreißen
\end{itemize}

\subsection{Individual Differences}\label{individual-differences}

Many decades of research in social cognition have taught us much about developmental milestones and at which average children possess certain social-cognitive abilities (for a review, see \citeproc{ref-rakoczy2022foundations}{Rakoczy, 2022}). For example, we have learned that children are around four to five year of age when they understand that other agents can have false beliefs that differ from the child's own belief and from reality (\citeproc{ref-wellman2001metaanalysis}{Wellman et al., 2001}).

Anyone who has ever interacted with children of a similar age must have noticed that they can still differ in their abilities: some children might be faster in learning how to walk, while others might be quicker in learning how to talk.

\emph{Individual differences} describe the features in which agents vary (``traits or other characteristics by which individuals may be distinguished from one another'', (\citeproc{ref-americanpsychologicalassociationindividual}{American Psychological Association, n.d.-a})). \emph{Inter}-individual differences focus on the variability \emph{between} children, while \emph{intra}-individual differences focus on the variability \emph{within} the same child. Asendorpf (\citeproc{ref-asendorpf1992stability}{1992}) have argued that inter-individual differences are often interpreted as trait-like characteristics (e.g.~personality) that are more or less stable over time, while intra-individual differences might capture state-like characteristics (e.g.~motivation) or development within a given construct. The present thesis and the associated publications focus on inter-individual differences, often abbreviated to individual differences.

A prominent framework that could explain differences between individuals is Vygotsky's sociocultural theory of cognitive development (\citeproc{ref-vygotsky1962thought}{Vygotsky, 1962}, \citeproc{ref-vygotsky1978mind}{1978}, \citeproc{ref-vygotsky1987thinking}{1987}). This theory emphasizes the role of social interaction in children's cognitive development and as such, learning as a social process. In other words, it is assumed that social interaction shapes social cognition.

In order to empirically test if social interaction shapes social cognition, we need to study children's daily environments and cognitive abilities on the individual level.

Studies focussing on individual differences aim to measure variation between children in a given social-cognitive ability and, potentially, study the co-development of certain cognitive abilities or link this variation to internal and external factors that influence its development. As such, individual differences studies act as a powerful tool to build and test developmental theories and the structure of social cognition (\citeproc{ref-happe2017structure}{Happé et al., 2017}). Additionally, individual differences studies are needed to design (clinical or educational) intervention programs. Capturing change before and after an intervention helps to evaluate its efficacy.

\begin{itemize}
\tightlist
\item
  also for modeling, understanding underlying processes?
\end{itemize}

Examples include studies conducted with twins (\citeproc{ref-hughes2005originsa}{Hughes et al., 2005}; \citeproc{ref-ronald2006individual}{Ronald et al., 2006}).

\begin{itemize}
\tightlist
\item
  example from comparative work: some researchers focus on what chimps in average do, others focus on the very end of the distribution: what are they capable of?
\item
  cross-cultural variation: not only average from small samples
\end{itemize}

False-belief understanding can predict children's pretense use, abilities for social interaction, peer popularity, which clearly underlines the relevance of individual differences in social-cognitive abilities for real-life outcomes (\citeproc{ref-wellman2013universal}{Wellman, 2013}).

\subsection{Cross-cultural Research}\label{cross-cultural-research}

Wellman (\citeproc{ref-wellman2013universal}{2013}) has argued that children's social-cognitive development is a particularly informative domain to assess in which ways people all over the world resemble an differ from one another.

\section{Gaze Following}\label{gaze-following}

\subsection{Terminology}\label{terminology-1}

\subsection{Developmental Milestones}\label{developmental-milestones-1}

\subsection{Individual Differences}\label{individual-differences-1}

\subsection{Cross-cultural Research}\label{cross-cultural-research-1}

\subsection{Research Gaps}\label{research-gaps}

\section{Methodological Considerations}\label{methodological-considerations}

\subsection{Terminology}\label{terminology-2}

In the context of this work, the terms \emph{variability}, \emph{variation} and \emph{variance} need further explanation. According to the APA Dictionary of Psychology, the definitions of variability and variation greatly resemble each other (variability: ``1. the quality of being subject to change or variation in behavior or emotion, 2. the degree to which members of a group or population differ from each other, as measured by statistics such as the range, standard deviation, and variance'' (\citeproc{ref-americanpsychologicalassociationvariability}{American Psychological Association, n.d.-d}); variation: ``1. the existence of qualitative differences in form, structure, behavior, and physiology among the individuals of a population, whether due to heredity or to environment. Both artificial selection and natural selection operate on variations among organisms, but only genetic variation is transmitted to offspring. 2. in statistics, the degree of variance or dispersion of values that is obtained for a specific variable'' (\citeproc{ref-americanpsychologicalassociationvariation}{American Psychological Association, n.d.-f})). Both definitions focus on differences between individuals that can be applied to the behavioral level. As in most of the developmental literature, the terms variability and variation will be used interchangeably in this work.

As the definitions of variability and variation indicate, the term \emph{variance} denotes a statistical measure for the degree to which group members differ from one another. It describes the dispersion of a distribution by calculating how far the values of a sample spread out from their average value. Consequently, a smaller variance means less differences between individuals in a given sample (a measure of the spread, or dispersion, of scores within a sample or population, whereby a small variance indicates highly similar scores, all close to the sample mean, and a large variance indicates more scores at a greater distance from the mean and possibly spread over a larger range'', (\citeproc{ref-americanpsychologicalassociationvariance}{American Psychological Association, n.d.-e})).

\emph{Validity} concerns the truthfulness of a measure or the degree to which a measure captures what it is designed to measure conceptually (``1. the characteristic of being founded on truth, accuracy, fact, or law. 2. the degree to which empirical evidence and theoretical rationales support the adequacy and appropriateness of conclusions drawn from some form of assessment {[}\ldots{]}'', (\citeproc{ref-americanpsychologicalassociationvalidity}{American Psychological Association, n.d.-c})). There are several different types of validity. In the context of this work, we focus on \emph{convergent} validity (\emph{i.e.,} the degree to which two measures that should be theoretically related are indeed related), \emph{predictive} validity (\emph{i.e.,} the degree to which a measure relates to other measures separated by a determined period (= predicting scores in the other measure)), and \emph{external} validity (\emph{i.e.,} the degree to which the results of a measure can be generalized beyond the original sample).

\emph{Reliability} refers to the degree to which a measure produces consistent, stable results when repeated again (``the trustworthiness or consistency of a measure, that is, the degree to which a test or other measurement instrument is free of random error, yielding the same results across multiple applications to the same sample'', (\citeproc{ref-americanpsychologicalassociationreliability}{American Psychological Association, n.d.-b}); ``). Two main types of reliability are often assessed when evaluating a given measure. \emph{Internal consistency} focuses on the relationship between items in the same measure to ensure they measure the same construct. \emph{(Test-) Retest reliability} assesses a measure's stability over time and estimates the relationship of the measurement scores on two different administration occasions (e.g., test day one followed by test day two three weeks later).

\subsection{Existing Social Cognition Measures}\label{existing-social-cognition-measures}

Unfortunately, there are not too many measures that were designed to capture individual differences in social cognition in children and according to Beaudoin et al. (\citeproc{ref-beaudoin2020systematic}{2020}), social cognition studies rarely report psychometrical information.

Reasons for such an insufficiency in individual differences measures might be the inherent challenges in methodological advancements. Developing a reliable and valid measure is time-consuming, and (due to some features of the publishing system) not overly rewarding.

An example is a longitudinal study by Sodian et al. (\citeproc{ref-sodian2016understanding}{2016}) who have investigated how goals, beliefs and desires relate to moral \hyperref[acronyms_ToM]{ToM} They reported that ``Contrary to expectations, neither goal encoding nor implicit false belief understanding was correlated with the Mo{[}ral{]} \hyperref[acronyms_ToM]{ToM} false belief question'' (p.~1227).

Two main interpretations come to mind when reading about these results. On one hand, the hypothesis that goal encoding and implicit FB understanding should be related to moral \hyperref[acronyms_ToM]{ToM} could be simply incorrect. On the other, more optimistic hand, the non-existent relationship could be explained by methodological issues: potentially, the applied measures could not capture variation between children. Interestingly, Sodian et al. (\citeproc{ref-sodian2016understanding}{2016}) stated two sentences later: ``{[}\ldots{]} verbal IQ was positively correlated with almost all the assessed variables'' (p.~1227). Why should we care about this relationship? The measures for verbal IQ was - by definition - designed to capture individual differences in linguistic intelligence. Possibly, the measures that we use for correlational analyses need to leave room for individual differences in order to find them.

What does it mean for a measure to be designed for individual differences? To drive the point home, let us imagine the opposite first. Many developmental study paradigms present children with two objects. For example, the Sally-Anne-task lets children witness how a ball gets moved from one box to a basked while another agent does or does not witness the change of location. Children then get asked where the agent believes the ball to be. Possible responses are (1) the box or (2) the basket. Similarly, in the traditional gaze following paradigm, an agent moves a toy into one of two boxes. Children themselves do not see where the toy was put but only see whether the agent looks to the (1) left- or (2) right-hand side. Common outcome measures focus on children's looking or reaching behavior toward the two boxes. What becomes apparent is that these common social cognition tasks rely on dichotomous measures: either the child passes or fails the task. The child's ability in question cannot be depicted on a continuous scale. A measure designed to capture individual differences would benefit from a continuous outcome measure, or, at least, more than one trial per child.

\subsection{The Need for a Paradigm Shift}\label{the-need-for-a-paradigm-shift}

Classical experimental paradigms often give little weigh to individual differences. When research questions focus on condition differences or the average age in which children pass a given social-cognitive task, variation between children is often handled by excluding outliers and calculating mean scores. As such, individual differences are regarded as noise or measurement error (\citeproc{ref-kidd2018individual}{Kidd et al., 2018}).

I will argue that we need to question the interpretation that variation between individuals equals noise and should be thrown away. Indeed, some research agendas rely on the assumption that meaningful individual differences exist. For example, correlational studies aim to quantify the degree of a relationship between two variables. They assess whether children that perform well in one task also perform well in another task. In order to assess such a relationship, an individual's ability must be measured accurately on both scales. Indeed, the correlation between two measures cannot be greater than the reliability of the individual measures. For example, let us assume we apply measure A with a known reliability estimate of .5 and measure B with a reliability estimate of .6 \ldots{}

Many research questions in developmental psychology are indeed individual differences questions.

Research focusing on the (dis)continuity of development must also rely on measures that can accurately track children's abilities across time. A textbook example is the question whether children develop some abilities similar to a tree that grows continuously over time or rather like a butterfly that develops in distinct stages (from caterpillar, to cocoon, to butterfly).

\begin{itemize}
\item
  An example of a gradual process developmental theory comes from Vygotsky. Theories within this realm assume that change occurs slowly and gradually.
\item
  An example of a stage-like developmental theory comes from Piaget who assumes that children \ldots{} Children are assumed to go through these qualitatively different stages in a set universal order.
\item
  First step is describing change, second step is explaining change! What changes over time and why? Process models.
\item
  Universality in development? Maybe absolute differences in test scores but similar patterns overall?
\item
  Interestingly, sometimes the same study paradigms are used for correlational research that were previously designed for experimental research focusing on group averages. It becomes apparent that this approach is
\end{itemize}

\subsection{A Potential Way Forward}\label{a-potential-way-forward}

That this endeavor is nevertheless worth it can be seen from a historical perspective on science. Often, only after new measurement techniques, new knowledge could be gained and theories could be advanced. In other words: newly designed methods led to new opportunities and research questions. {[}todo example, maybe Teleskop{]}

\begin{itemize}
\item
  correlations only as large as the least reliable measure; we correlate lots but rarely know the psychometrics of the tasks
\item
  potentiall issues:

  \begin{itemize}
  \tightlist
  \item
    group-level averages
  \item
    small sample sizes
  \item
    low trial numbers
  \item
    dichotomous measures
  \item
    missing or unsatisfactory psychometric properties
  \end{itemize}
\item
  Poor measurement on an individual level might conceal relationships between different aspects of cognition and may obscure developmental change
\item
  Wishlist:

  \begin{itemize}
  \tightlist
  \item
    Objective, standardized measures (compared to anecdotal evidence, diaries)
  \item
    Parametric measures (= continuous instead of dichotomous. Avoid floor \& ceiling effects) (Schaafsma, Pfaff, Spunt, and Adolphs, 2015)
  \item
    Satisfactory reliability estimates (Beaudoin et al., 2020; Hughes \& Devine, 2015)
  \item
    Induce variation across individuals and age groups (Repacholi, 2003)
  \end{itemize}
\item
  Behavior measured should represent behavior in real world (Repacholi, 2003)
\item
  many seemingly unrelated issues: replicability crisis, no correlations when theoretically expected, overreliance on Global North samples
\item
  how can we address these issues? maybe for many similar solution, namely robust methods
\item
  what makes a method robust? vali \& reli
\item
  but also: to move science forward, collaborative endeveaours are needed. Open science, share the task, so that others can reproduce results and test theories \& generalizability
\item
  we need new measures to capture individual differences in social cognition
\item
  this thesis does so for a fundamental social-cognitive ability: gaze following
\item
  paradigmen shift: here example of one construct how this might be possible
\item
  capture ind diff reliably. check process behind it and related constructs. see whether this is universally applicable
\end{itemize}

\chapter{General Approach}\label{general-approach}

\section{Aims of this Dissertation}\label{aims-of-this-dissertation}

The overarching aim of this dissertation is to deepen our understanding of variability in a core social-cognitive ability, namely gaze following. In the {[}Introduction{]}, I have highlighted calls for better (\emph{i.e.,} valid, reliable, generalizable) social cognition measures. The projects included in this dissertation follow this call --- and try to demonstrate that, even if demanding, methodological advancement is feasible. By focusing on an individual differences measure, this dissertation aims to question the conception that variation between participants equals noise and should be disregarded. By applying the newly designed gaze following measure, I hope to show-case how individual differences research on social cognition helps us to answer fundamental questions about the development, cognitive processes and universality of one of the most extraordinary human abilities. In the following, I will outline the role each of the four included studies play in reaching these overall objectives.

\subsection{Aims of Study I}\label{aims-of-study-i}

Much of the existing social cognition research relies on study designs that were created for group-level analyses. Yet, many research questions focus on individual differences: for example, do siblings have a positive influence on children's social-cognitive development (\citeproc{ref-devine2018family}{Devine \& Hughes, 2018}; \citeproc{ref-perner1994theory}{Perner et al., 1994}; \citeproc{ref-peterson2000kindred}{Peterson, 2000}; \citeproc{ref-stewart1983sibling}{Stewart, 1983})? How are gaze following and language abilities related (\citeproc{ref-brooks2005development}{Brooks \& Meltzoff, 2005}, \citeproc{ref-brooks2015connecting}{2015}; \citeproc{ref-macdonald2013eye}{Macdonald \& Tatler, 2013}; \citeproc{ref-okumura2017individual}{Okumura et al., 2017})? Can a child's perspective-taking ability be trained (\citeproc{ref-lecce2014promoting}{Lecce et al., 2014}; \citeproc{ref-ruffman2018variety}{Ruffman et al., 2018}; \citeproc{ref-zhang2021worse}{Zhang et al., 2021})?

In order to relate cognitive abilities and external factors to each other, social cognition tasks need to be able to measure differences between children (or within children over a certain period of time). Tasks that are designed to do so remain rare. This study aimed at approaching this issue and described the development and validation of our new gaze following task, the TANGO (Task for Assessing iNdividual differences in Gaze understanding --- Open). We have designed a playful study in which participants are asked to locate a balloon. The twist is that the participants cannot see the balloon itself. However, they see how another agent gazes out of a window, into the direction of the balloon. Participants have to follow the agent's gaze to locate the balloon. Their precision in doing so is measured continuously, namely in the distance between the actual target's location and their assumed location.

\hyperref[studyI]{Study I} consisted of three main parts. In the first part, we tested the task's usability in a sample of 3- to 5-year-old children and adults. We assessed whether the TANGO captured individual differences and compared data collected remotely vs.~in-person (for the child sample). In the second part, we assessed whether these individual differences were reliable. We repeated data collection after two weeks and estimated retest-reliability and internal consistency. In the third part, we assessed the TANGO's validity. For its construct validity, we explored the relationship between the TANGO and variables of children's daily social environment. For its predictive validity, we examined the TANGO's relationship with a language measure six months later. Together, these findings thoroughly assessed the TANGO as a new task to measure individual differences in gaze following.

\subsection{Aims of Study II}\label{aims-of-study-ii}

\hyperref[studyII]{Study II} adopted the TANGO to examine three research questions related to individual differences: (1) how does gaze following develop across the lifespan, (2) how do people process gaze cues, and (3) how does gaze following relate to other (social-) cognitive abilities.

To answer the first question, we remotely collected data from 3- to 80-year-olds and examined changes in their ability to follow gaze. While previous gaze following research often focused on the youngest age in which children can follow gaze (\citeproc{ref-dentremont1997demonstration}{D'Entremont et al., 1997}; e.g., \citeproc{ref-dentremont2000perceptual}{D'Entremont, 2000}), we focused on the development (in precision) after this initial milestone.

Second, we developed a computational cognitive model to explain the underlying processes of how people process gaze. We hypothesized that participants observe the pupil location within the eye and calculate gaze vectors that point toward the attentional focus. As this process was likely to be noisy, individual differences were assumed in the participant's uncertainty surrounding this gaze vector. We judged the evidence for this gaze model by comparing it to two alternative explanations of the data: a center bias and a random guessing model. Due to its underlying geometrical features, the gaze model predicts that uncertainties around the gaze vector increase the further to the side an agent looks. We checked for this signature pattern in the data by assessing the participants' imprecision levels in each target bin.

Third, we tested whether an individual's precision in gaze following relates to other (social-) cognitive abilities. Consistent with our gaze model, we assessed whether gaze following was related to non-social vector following. To test this claim, we designed a new vector following task that matched the TANGO as closely as possible but within a non-social context. The vector following task relied on the concept of magnetism and asked children to locate a magnet based on the location of a gearwheel. Furthermore, we applied four tasks out of a \hyperref[acronyms_ToM]{ToM} task battery (\citeproc{ref-wellman2001metaanalysis}{Wellman et al., 2001}) and two perspective-taking tasks (\citeproc{ref-flavell1981younga}{Flavell, Everett, et al., 1981}; \citeproc{ref-flavell1981development}{Flavell, Flavell, et al., 1981}). These comparisons aimed at disentangling the components that make up gaze following.

\subsection{Aims of Study III}\label{aims-of-study-iii}

The studies presented above relied on child samples from the Global North (Leipzig, Germany). Developmental theories often assume that central features of social cognition hold universally true across human cultures (\citeproc{ref-wellman2013universal}{Wellman, 2013}). Study III aimed at evaluating this claim: To broaden our perspective, we collected data of children's gaze following abilities from 17 diverse communities spanning five different continents. To our knowledge, this constitutes the largest international sample on children's gaze following.

First, we examined the developmental trajectory of gaze following across the communities. Second, we explored whether absolute differences in gaze following imprecision related to variables concerning the data collection mode and/or children's daily environments. Third, we checked whether the signature pattern predicted by our gaze model in Study II could be found in all communities studied. With this approach, we aimed to find out whether the process of children follow gaze is similar worldwide.

\subsection{Aims of Study IV}\label{aims-of-study-iv}

\hyperref[studyIV]{Study IV} strengthens the basis for using the newly developed gaze following task for cross-cultural research. The paper had three main aims.

First, we described the process of designing the cross-cultural gaze following task (TANGO---CC). We aimed to provide a roadmap for other researchers how to pragmatically develop a new task that can be used for studying diverse communities.

Second, we provided a tutorial for how other researchers can use and customize the TANGO---CC according to their needs. This section aimed to highlight the components of the task that are constant and those that are variable and can be adapted to diverse languages and communities.

Third (and most importantly), we aimed at examining the TANGO---CC's reliability across diverse communities. While the psychometric properties of the TANGO were previously tested in one setting (Leipzig, Germany; \hyperref[studyI]{Study I}), we cannot simply generalize these findings and assume the task's validity and reliability based on a mono-cultural sample. For this purpose, we further analyzed the data set from \hyperref[studyIII]{Study III}. We assessed performance across trial types, within- and between-community variation, and internal consistency measures in all of the 17 communities studied.

\section{Research Practices}\label{research-practices}

Above, I have highlighted the theoretical approaches and aims of the publications included in this dissertation. All of these studies build upon the same underlying testing infrastructure. As this dissertation has a strong focus on methodological advances, I want to outline our main technological decisions within the process of designing this new testing infrastructure. In-depth information regarding each specific study can be found in the corresponding manuscripts in {[}Appendix A{]}.

\subsection{Accessibility of Study Materials}\label{accessibility-of-study-materials}

Elson et al. (\citeproc{ref-elson2023psychological}{2023}) figuratively stated that ``psychological measures aren't toothbrushes'' (p.1). The authors argue that psychological tasks are often used only a handful of times --- by the same researcher(s) who designed the measure. However, this reluctance to share own measures and re-use measures of other researchers hinders cumulative science. Elson and colleagues, therefore, call for more transparency and standardization in the psychological tasks we implement and apply. We took this call seriously and strongly adhered to Open Science principles in the construction of our new standardized gaze following task.

The projects included in this dissertation pre-registered studies with their planned sample sizes, study procedures and planned analyses for hypothesis testing before data collection, as well as shared the study materials, analysis scripts and data sets after data collection. This dissertation and its included publications have published all related information in repositories on GitHub (\url{https://github.com}) and the Open Science Framework (\url{https://osf.io}; links leading to specific projects can be found in the respective manuscripts in the {[}Appendix A{]}). Articles were written in R Markdown which allows the user to combine continuous text and into one coherent document. Tracking analysis scripts and the manuscript texts in a version-controlled manner lets interested researchers retrace decisions and changes along the way. The newly implemented task(s) are openly available to other researchers to use and share as they wish. The source code of the website is also available via GitHub. By cloning the corresponding repository, motivated researchers can modify and further develop the task to their needs.

As the final project, \hyperref[studyIV]{Study IV} incorporates refinements that have further optimized the usability of the gaze following task. While adapting the source code for cross-cultural stimulus presentation, I have refactored the code base (\emph{i.e.}, improved the internal structure, readability, and maintainability of the software without altering its behavior or functionality). I have followed added comments to the code base using JSDoc (\url{https://jsdoc.app}), a markup language developed to document and annotate JavaScript source code files. Furthermore, I have collected frequently asked questions and created a manual for other researchers using the task (\url{https://ccp-odc.eva.mpg.de/tango-cc/manual.html}). These changes compared to the first version of the gaze following task hopefully serve to improve its sustainability over time.

\subsection{Presentation Mode}\label{presentation-mode}

Before designing the gaze following task itself, we had to choose a presentation mode for the experimental stimuli. Many traditional developmental studies have designed interactive tasks for children in which they interact with the experimenter or their caregiver. This makes standardization difficult as experimenters have individual communication styles, might accidentally alter the instruction script due to memory gaps, or need replacement at some point (e.g., when temporary work contracts end). In addition, this approach makes scaling data collection difficult as training experimenters requires substantial temporal and financial resources. Since studies investigating individual differences rely on large sample sizes, we decided to take another route. We designed an online testing platform that hosts studies in all common web browsers. Please note that even though the task is programmed as a website, it does not require a working WiFi connection (for more information, see further down).

The design as an online study has several advantages. First, it allows for maximally standardized procedures as audio instructions are pre-recorded and automated. Second, data does not need to be entered manually or coded from video. Instead, the website stores response times and clicking behavior automatically. This greatly reduces potential biases and rater errors. Third, the presentation as a website allows for in-person and remote study participation. Families can choose whether they want to participate as before by coming to the research lab or kindergartens, or whether they want to participate from the comfort of their own home. In this case, no appointments or rooms need to be scheduled. The families receive a personalized invitation link and can access the online study whenever and wherever it suits them best. This increases the inclusiveness for families with busy schedules, longer traveling distances to the lab or children who might be shy to interact with strangers. While studies conducted in a research lab might be prone to over-study families with certain characteristics (e.g., families living closer by, having a higher SES or more free time), online studies might make study participation more accessible for a wider range of families --- which, in turn, would increase the representativeness of the sample.

We have consciously chosen to implement an active behavioral measure as we wanted to circumvent issues in interpretability of implicit measures. For example, it is unclear whether looking time measures capture children's level of surprise, attention, memory formation or other cognitive processes, and whether they show convergent validity (\citeproc{ref-aslin2007look}{Aslin, 2007}; \citeproc{ref-dorrenberg2018how}{Dörrenberg et al., 2018}). While our outcome response does not provide a direct measure of underlying cognitive processes, it focuses on children's active behavior which might translate more directly into their actions in daily life.

The task itself is presented as an animated, interactive picture book. Audio instructions guide children through the task and provide a description of the presented events. Trial types build up on each other and familiarize children with the response format.

\subsection{Programming Framework}\label{programming-framework}

The online study was programmed in JavaScript (functionality of the website), HTML (content of the website), CSS (styling of the website) and PHP (server-side interaction for down-/ uploading data). We decided to refrain from using web development frameworks like Next.js or React. As the web development is a very fast moving domain with frequently appearing technology advances, projects relying on new frameworks often require regular monitoring and updates. Our goal was to create a durable code base that avoids breaking changes through large updates. Additionally, we aimed at keeping the code base as easily accessible as possible so that also researchers with less programming experience can read through and adjust the code. We reasoned that this could be better achieved by relying on the so-called ``Vanilla JavaScript'' without additional frameworks. An exception to this was an animation library we used. The GreenSock Animation Platform (GSAP; \url{https://gsap.com}) is a JavaScript library for industry-standard, high-performance web animations. However, GSAP is widely spread, well documented, and easy to read so including this library rather increases user friendliness. We used one further tool, namely a module bundler for JavaScript. Parcel (\url{https://parceljs.org}), and in later projects, webpack (\url{https://webpack.js.org}) tracked modules and dependencies, optimized website performance and enhanced the developer experience (e.g., by enabling a development server and hot reloading). However, please note that these bundlers only alter the source code on the sever-side --- the functions, etc., defined in Vanilla JavaScript remain in their human-readable format.

During website development, we made use of the development server made available by the module bundler. This means that we could test the website on our local computer and only uploaded content to the web server once the website was fully implemented. The bundler builds a so-called ``dist'' folder that contains all website content in a compressed and optimized format. This folder is essential for both online and offline data collection. For online data collection, the dist folder gets uploaded to the web server and the study can be accessed under its corresponding URL. For offline data collection, a live server can be installed and hosts the dist folder, so that the website can be displayed on a computer without any internet connection. This is especially practical for data collection in settings with unreliable WiFi. Potentially, this might help for collecting data in communities living in remote villages, but also for data collection with children in (German) kindergartens as these institutions often do not have public internet connections or telephone reception.

\subsection{Data Processing}\label{data-processing}

After a child completed the task, the website automatically saves all responses and response times. The website stores all click data and, if desired, some additional device information (e.g., touch screen yes/no, what type and version of web browser) in a text format. During initial development and the first published article (see \hyperref[studyI]{Study I}), we used JSON (JavaScript Object Notation) files which are commonly used for transmitting data in web applications. In later stages of the project (see \hyperref[studyIII]{Study III}), we instead used CSV (Comma-Separated Values) files as these are more commonly used among psychology researchers. In these text files, each trial is represented by one row of data. Columns depict the collected variables, for example, subject id (anonymized), trial type, and response. If requested, a webcam recording can be captured that films the study participation via the front camera. We implemented this feature especially for remote data collection. After the consent of the caregivers, we could use the webcam recordings to ensure that indeed the children themselves carried out the task and did not get any help from their caregivers. Depending on the choice of the experimenter and the data collection setup, data will be uploaded to secure servers located within the Max Planck Institute for Evolutionary Anthropology, or will be downloaded to the testing device (e.g., tablet or computer).

Data pre-processing was greatly simplified as the study data does not require any further response coding. This saves time and effort by experimenters and reduces potential issues introduced by rater errors. Individual subject files were merged into one final data set using the programming language R. Data visualization and statistical analyses were conducted within the same program.

\chapter{General Discussion}\label{discussion}

\section{Summary of Results}\label{summary-of-results}

\subsection{Results of Study I}\label{resultsI}

The task design

\subsection{Results of Study II}\label{resultsII}

\subsection{Results of Study III}\label{resultsIII}

\subsection{Results of Study IV}\label{resultsIV}

\section{Research Contributions}\label{research-contributions}

\begin{itemize}
\tightlist
\item
  also include orev \& previc as examples how the platform idea is useful tool (gleiche Plattform, Prinzipien der Entwicklung)
\end{itemize}

\section{Limitations}\label{limitations}

\subsection{The task itself}\label{the-task-itself}

\begin{itemize}
\tightlist
\item
  social enough? still on screen, only observing
\item
  idea for first social interaction, where participant guides where agent looks. more like a second-person perspective
\item
  change in framing: gaze understanding to gaze following. why?
\end{itemize}

Developmental studies often report high number of drop-out rates, for example because of children's fuzziness (todo: example). Administering the The TANGO(---CC) has proven very helpful in this regard: across communities, children tended to complete the task with ease and even asked for a repetition. Drop-out rates could thus be greatly reduced. As the data collection procedures varied slightly from community to community, we cannot report exact numbers of drop-outs but rely on the experience of experimenters.

\subsection{Modeling}\label{modeling}

\begin{itemize}
\tightlist
\item
  why is gafo still social if vector?
\item
  why model needed? cc signature pattern, other interpretation than raw scores, in future incorporate more complex mental states
\item
  why new? compared to others, only reliance on eyes not head. plus ind diff
\end{itemize}

\subsection{Cross-cultural work}\label{cross-cultural-work}

\begin{itemize}
\tightlist
\item
  do not want to signal cc is easy or we propose parachute science
\item
  technology effect: yes but still ind diff in communities where touchscreen = 1
\item
  more helpful to check relative imprecision, not absolute values within or between communities
\item
  idea for touchscreen familiarization: bubble wrap game
\item
  (\citeproc{ref-wellman2013universal}{Wellman, 2013}) how can social cognition be universal and yet vary? framework understanding
\item
  importantly, focusing on signature patterns helps us to study underlying cognitive processes, which raw absolutes score could not have revealed
\item
  don't just compare urban global north to rural global south, which is common pitfall in CC work (Barrett 2020)
\end{itemize}

\section{Outlook}\label{outlook}

Applying the tool: Chapter on how Gafo can be extended to answer new research questions.

\subsection{Variations of TANGO}\label{variations-of-tango}

\begin{itemize}
\tightlist
\item
  gafo alien to answer why we need to integrate information from 2 eyes. probably important for side locations in tango?
\item
  circle like clock to circumvent precision with distance diffusion
\item
  different starting directions to avoid center bias
\end{itemize}

\subsection{Open research questions}\label{open-research-questions}

\begin{itemize}
\tightlist
\item
  but also: all methods sanity checks so that now theory testing can start
\item
  do social interactions shape social cognition abilities? we need good measures on both ends, machine learning as promising avenue
\item
  back to uniquely human: what about great apes? how would they perform in gafo?
\item
  now that we can study vector following, we can also go down this path: is vector estimation needed in other social cognitive abilities? what about for example action prediction, reaching at sth?
\end{itemize}

\section{Conclusion}\label{conclusion}

\backmatter

\chapter{References}\label{references}

\phantomsection\label{refs}
\begin{CSLReferences}{1}{0}
\bibitem[\citeproctext]{ref-americanpsychologicalassociationindividual}
American Psychological Association. (n.d.-a). Individual {Differences}. In \emph{APA Dictionary of Psychology}. Retrieved June 26, 2024.

\bibitem[\citeproctext]{ref-americanpsychologicalassociationreliability}
American Psychological Association. (n.d.-b). Reliability. In \emph{APA Dictionary of Psychology}. Retrieved June 26, 2024.

\bibitem[\citeproctext]{ref-americanpsychologicalassociationvalidity}
American Psychological Association. (n.d.-c). Validity. In \emph{APA Dictionary of Psychology}. Retrieved June 26, 2024.

\bibitem[\citeproctext]{ref-americanpsychologicalassociationvariability}
American Psychological Association. (n.d.-d). Variability. In \emph{APA Dictionary of Psychology}. Retrieved June 26, 2024.

\bibitem[\citeproctext]{ref-americanpsychologicalassociationvariance}
American Psychological Association. (n.d.-e). Variance. In \emph{APA Dictionary of Psychology}. Retrieved June 26, 2024.

\bibitem[\citeproctext]{ref-americanpsychologicalassociationvariation}
American Psychological Association. (n.d.-f). Variation. In \emph{APA Dictionary of Psychology}. Retrieved June 26, 2024.

\bibitem[\citeproctext]{ref-asendorpf1992stability}
Asendorpf, J. B. (1992). Beyond stability: {Predicting} inter-individual differences in intra-individual change. \emph{European Journal of Personality}, \emph{6}(2), 103--117. \url{https://doi.org/10.1002/per.2410060204}

\bibitem[\citeproctext]{ref-aslin2007look}
Aslin, R. N. (2007). What's in a look? \emph{Developmental Science}, \emph{10}(1), 48--53. \url{https://doi.org/10.1111/j.1467-7687.2007.00563.x}

\bibitem[\citeproctext]{ref-beaudoin2020systematic}
Beaudoin, C., Leblanc, É., Gagner, C., \& Beauchamp, M. H. (2020). Systematic {Review} and {Inventory} of {Theory} of {Mind Measures} for {Young Children}. \emph{Frontiers in Psychology}, \emph{10}, 2905. \url{https://doi.org/10.3389/fpsyg.2019.02905}

\bibitem[\citeproctext]{ref-bowlby1958nature}
Bowlby, J. (1958). The nature of the child's tie to his mother. \emph{The International Journal of Psychoanalysis}, \emph{39}, 350--373.

\bibitem[\citeproctext]{ref-brooks2005development}
Brooks, R., \& Meltzoff, A. N. (2005). The development of gaze following and its relation to language. \emph{Developmental Science}, \emph{8}(6), 535--543. \url{https://doi.org/10.1111/j.1467-7687.2005.00445.x}

\bibitem[\citeproctext]{ref-brooks2015connecting}
Brooks, R., \& Meltzoff, A. N. (2015). Connecting the dots from infancy to childhood: {A} longitudinal study connecting gaze following, language, and explicit theory of mind. \emph{Journal of Experimental Child Psychology}, \emph{130}, 67--78. \url{https://doi.org/10.1016/j.jecp.2014.09.010}

\bibitem[\citeproctext]{ref-csibra2009natural}
Csibra, G., \& Gergely, G. (2009). Natural pedagogy. \emph{Trends in Cognitive Sciences}, \emph{13}(4), 148--153. \url{https://doi.org/10.1016/j.tics.2009.01.005}

\bibitem[\citeproctext]{ref-dentremont2000perceptual}
D'Entremont, B. (2000). A perceptual--attentional explanation of gaze following in 3- and 6-month-olds. \emph{Developmental Science}, \emph{3}(3), 302--311. \url{https://doi.org/10.1111/1467-7687.00124}

\bibitem[\citeproctext]{ref-dentremont1997demonstration}
D'Entremont, B., Hains, S. M. J., \& Muir, D. W. (1997). A demonstration of gaze following in 3- to 6-month-olds. \emph{Infant Behavior and Development}, \emph{20}(4), 569--572. \url{https://doi.org/10.1016/S0163-6383(97)90048-5}

\bibitem[\citeproctext]{ref-decety2020preface}
Decety, J. (2020). Preface. In J. Decety (Ed.), \emph{The {Social Brain}: {A Developmental Perspective}} (pp. ix--xii). The MIT Press.

\bibitem[\citeproctext]{ref-devine2018family}
Devine, R. T., \& Hughes, C. (2018). Family {Correlates} of {False Belief Understanding} in {Early Childhood}: {A Meta-Analysis}. \emph{Child Development}, \emph{89}(3), 971--987. \url{https://doi.org/10.1111/cdev.12682}

\bibitem[\citeproctext]{ref-dorrenberg2018how}
Dörrenberg, S., Rakoczy, H., \& Liszkowski, U. (2018). How (not) to measure infant {Theory} of {Mind}: {Testing} the replicability and validity of four non-verbal measures. \emph{Cognitive Development}, \emph{46}, 12--30. \url{https://doi.org/10.1016/j.cogdev.2018.01.001}

\bibitem[\citeproctext]{ref-elson2023psychological}
Elson, M., Hussey, I., Alsalti, T., \& Arslan, R. C. (2023). Psychological measures aren't toothbrushes. \emph{Communications Psychology}, \emph{1}(1), 1--4. \url{https://doi.org/10.1038/s44271-023-00026-9}

\bibitem[\citeproctext]{ref-flavell1981younga}
Flavell, J. H., Everett, B. A., Croft, K., \& Flavell, E. R. (1981). Young children's knowledge about visual perception: {Further} evidence for the {Level} 1--{Level} 2 distinction. \emph{Developmental Psychology}, \emph{17}(1), 99--103. \url{https://doi.org/10.1037/0012-1649.17.1.99}

\bibitem[\citeproctext]{ref-flavell1981development}
Flavell, J. H., Flavell, E. R., Green, F. L., \& Wilcox, S. A. (1981). The {Development} of {Three Spatial Perspective-Taking Rules}. \emph{Child Development}, \emph{52}(1), 356--358. \url{https://doi.org/10.2307/1129250}

\bibitem[\citeproctext]{ref-glynn2016social}
Glynn, M. A., \& Watkiss, L. (2016). Social {Cognition}. In M. Augier \& D. J. Teece (Eds.), \emph{The {Palgrave Encyclopedia} of {Strategic Management}} (pp. 1--4). Palgrave Macmillan UK. \url{https://doi.org/10.1057/978-1-349-94848-2_614-1}

\bibitem[\citeproctext]{ref-grossmann2022becoming}
Grossmann, T. (2022). Becoming uniquely human? {Comparing} chimpanzee to human infancy. \emph{Developmental Science}, \emph{25}(1), e13142. \url{https://doi.org/10.1111/desc.13142}

\bibitem[\citeproctext]{ref-happe2017structure}
Happé, F., Cook, J. L., \& Bird, G. (2017). The {Structure} of {Social Cognition}: {In}(ter)dependence of {Sociocognitive Processes}. \emph{Annual Review of Psychology}, \emph{68}(1), 243--267. \url{https://doi.org/10.1146/annurev-psych-010416-044046}

\bibitem[\citeproctext]{ref-henrich2015secreta}
Henrich, J. (2015). \emph{The {Secret} of {Our Success}: {How Culture Is Driving Human Evolution}, {Domesticating Our Species}, and {Making Us Smarter}}. Princeton University Press. \url{https://doi.org/10.1515/9781400873296}

\bibitem[\citeproctext]{ref-herrmann2007humansa}
Herrmann, E., Call, J., Hernàndez-Lloreda, M. V., Hare, B., \& Tomasello, M. (2007). Humans have evolved specialized skills of social cognition: The cultural intelligence hypothesis. \emph{Science (New York, N.Y.)}, \emph{317}(5843), 1360--1366. \url{https://doi.org/10.1126/science.1146282}

\bibitem[\citeproctext]{ref-hughes2005originsa}
Hughes, C., Jaffee, S. R., Happé, F., Taylor, A., Caspi, A., \& Moffitt, T. E. (2005). Origins of individual differences in theory of mind: From nature to nurture? \emph{Child Development}, \emph{76}(2), 356--370. \url{https://doi.org/10.1111/j.1467-8624.2005.00850.x}

\bibitem[\citeproctext]{ref-kidd2018individual}
Kidd, E., Donnelly, S., \& Christiansen, M. H. (2018). Individual {Differences} in {Language Acquisition} and {Processing}. \emph{Trends in Cognitive Sciences}, \emph{22}(2), 154--169. \url{https://doi.org/10.1016/j.tics.2017.11.006}

\bibitem[\citeproctext]{ref-lecce2014promoting}
Lecce, S., Bianco, F., Devine, R. T., Hughes, C., \& Banerjee, R. (2014). Promoting theory of mind during middle childhood: {A} training program. \emph{Journal of Experimental Child Psychology}, \emph{126}, 52--67. \url{https://doi.org/10.1016/j.jecp.2014.03.002}

\bibitem[\citeproctext]{ref-macdonald2013eye}
Macdonald, R. G., \& Tatler, B. W. (2013). Do as eye say: {Gaze} cueing and language in a real-world social interaction. \emph{Journal of Vision}, \emph{13}(4), 6--6. \url{https://doi.org/10.1167/13.4.6}

\bibitem[\citeproctext]{ref-moskowitz2013social}
Moskowitz, G. B. (2013). Social {Cognition}. \emph{Obo in Psychology}. \url{https://doi.org/10.1093/obo/9780199828340-0099}

\bibitem[\citeproctext]{ref-okumura2017individual}
Okumura, Y., Kanakogi, Y., Kobayashi, T., \& Itakura, S. (2017). Individual differences in object-processing explain the relationship between early gaze-following and later language development. \emph{Cognition}, \emph{166}, 418--424. \url{https://doi.org/10.1016/j.cognition.2017.06.005}

\bibitem[\citeproctext]{ref-ostrom1994foreword}
Ostrom, T. M. (1994). Foreword. In S. R. Wyer \& T. K. Srull (Eds.), \emph{Handbook of {Social Cognition}: {Applications}.} (2nd ed., Vol. 2, pp. vii--xii). Lawrence Erlbaum Associates.

\bibitem[\citeproctext]{ref-perner1994theory}
Perner, J., Ruffman, T., \& Leekam, S. R. (1994). Theory of {Mind Is Contagious}: {You Catch It} from {Your Sibs}. \emph{Child Development}, \emph{65}(4), 1228--1238. \url{https://doi.org/10.2307/1131316}

\bibitem[\citeproctext]{ref-peterson2000kindred}
Peterson. (2000). Kindred spirits: {Influences} of siblings' perspectives on theory of mind. \emph{Cognitive Development}, \emph{15}(4), 435--455. \url{https://doi.org/10.1016/S0885-2014(01)00040-5}

\bibitem[\citeproctext]{ref-pontius2007initial}
Pontius, J. U., Mullikin, J. C., Smith, D. R., Team, A. S., Lindblad-Toh, K., Gnerre, S., Clamp, M., Chang, J., Stephens, R., Neelam, B., Volfovsky, N., Schäffer, A. A., Agarwala, R., Narfström, K., Murphy, W. J., Giger, U., Roca, A. L., Antunes, A., Menotti-Raymond, M., \ldots{} O'Brien, S. J. (2007). Initial sequence and comparative analysis of the cat genome. \emph{Genome Research}, \emph{17}(11), 1675--1689. \url{https://doi.org/10.1101/gr.6380007}

\bibitem[\citeproctext]{ref-prein2023tango}
Prein, J. C., Kalinke, S., Haun, D. B. M., \& Bohn, M. (2023). {TANGO}: {A} reliable, open-source, browser-based task to assess individual differences in gaze understanding in 3 to 5-year-old children and adults. \emph{Behavior Research Methods}, \emph{56}(3), 2469--2485. \url{https://doi.org/10.3758/s13428-023-02159-5}

\bibitem[\citeproctext]{ref-quesque2024defining}
Quesque, F., Apperly, I., Baillargeon, R., Baron-Cohen, S., Becchio, C., Bekkering, H., Bernstein, D., Bertoux, M., Bird, G., Bukowski, H., Burgmer, P., Carruthers, P., Catmur, C., Dziobek, I., Epley, N., Erle, T. M., Frith, C., Frith, U., Galang, C. M., \ldots{} Brass, M. (2024). Defining key concepts for mental state attribution. \emph{Communications Psychology}, \emph{2}(1), 1--5. \url{https://doi.org/10.1038/s44271-024-00077-6}

\bibitem[\citeproctext]{ref-quesque2020theoryofmind}
Quesque, F., \& Rossetti, Y. (2020). What {Do Theory-of-Mind Tasks Actually Measure}? {Theory} and {Practice}. \emph{Perspectives on Psychological Science}, \emph{15}(2), 384--396. \url{https://doi.org/10.1177/1745691619896607}

\bibitem[\citeproctext]{ref-rakoczy2022foundations}
Rakoczy, H. (2022). Foundations of theory of mind and its development in early childhood. \emph{Nature Reviews Psychology}, \emph{1}(4), 223--235. \url{https://doi.org/10.1038/s44159-022-00037-z}

\bibitem[\citeproctext]{ref-ronald2006individual}
Ronald, A., Viding, E., Happé, F., \& Plomin, R. (2006). Individual differences in theory of mind ability in middle childhood and links with verbal ability and autistic traits: {A} twin study. \emph{Social Neuroscience}, \emph{1}(3-4), 412--425. \url{https://doi.org/10.1080/17470910601068088}

\bibitem[\citeproctext]{ref-ruffman2018variety}
Ruffman, T., Puri, A., Galloway, O., Su, J., \& Taumoepeau, M. (2018). Variety in parental use of {``want''} relates to subsequent growth in children's theory of mind. \emph{Developmental Psychology}, \emph{54}(4), 677. \url{https://doi.org/10.1037/dev0000459}

\bibitem[\citeproctext]{ref-sodian2016understanding}
Sodian, B., Licata, M., Kristen-Antonow, S., Paulus, M., Killen, M., \& Woodward, A. (2016). Understanding of {Goals}, {Beliefs}, and {Desires Predicts Morally Relevant Theory} of {Mind}: {A Longitudinal Investigation}. \emph{Child Development}, \emph{87}(4), 1221--1232. \url{https://doi.org/10.1111/cdev.12533}

\bibitem[\citeproctext]{ref-stewart1983sibling}
Stewart, R. B. (1983). Sibling {Interaction}: {The Role} of the {Older Child} as {Teacher} for the {Younger}. \emph{Merrill-Palmer Quarterly}, \emph{29}(1), 47--68. \url{https://www.jstor.org/stable/23086191}

\bibitem[\citeproctext]{ref-tomasello2019becoming}
Tomasello, M. (2019). \emph{Becoming {Human}: {A Theory} of {Ontogeny}}. Harvard University Press. \url{https://doi.org/10.4159/9780674988651}

\bibitem[\citeproctext]{ref-tomasello2020adaptivea}
Tomasello, M. (2020). The adaptive origins of uniquely human sociality. \emph{Philosophical Transactions of the Royal Society B: Biological Sciences}, \emph{375}(1803), 20190493. \url{https://doi.org/10.1098/rstb.2019.0493}

\bibitem[\citeproctext]{ref-vygotsky1962thought}
Vygotsky, L. S. (1962). \emph{Thought and language} (E. Hanfmann \& G. Vakar, Eds.; pp. xxi, 168). MIT Press. \url{https://doi.org/10.1037/11193-000}

\bibitem[\citeproctext]{ref-vygotsky1978mind}
Vygotsky, L. S. (1978). \emph{Mind in {Society}: {Development} of {Higher Psychological Processes}} (M. Cole, V. Jolm-Steiner, S. Scribner, \& E. Souberman, Eds.). Harvard University Press. \url{https://doi.org/10.2307/j.ctvjf9vz4}

\bibitem[\citeproctext]{ref-vygotsky1987thinking}
Vygotsky, L. S. (1987). Thinking and speech. In R. W. Rieber \& A. S. Carton (Eds.), \emph{The collected works of {L}.{S}. {Vygotsky}, {Volume} 1: {Problems} of general psychology} (Vol. 11, pp. 39--285). Plenum Press. (Original work published 1934).

\bibitem[\citeproctext]{ref-waterson2005initial}
Waterson, R. H., Lander, E. S., Wilson, R. K., \& The Chimpanzee Sequencing and Analysis Consortium. (2005). Initial sequence of the chimpanzee genome and comparison with the human genome. \emph{Nature}, \emph{437}(7055), 69--87. \url{https://doi.org/10.1038/nature04072}

\bibitem[\citeproctext]{ref-wellman2013universal}
Wellman, H. M. (2013). Universal {Social Cognition}: {Childhood Theory} of {Mind}. In M. R. Banaji \& S. A. Gelman (Eds.), \emph{Navigating the {Social World}: {What Infants}, {Children}, and {Other Species Can Teach Us}} (p. 0). Oxford University Press. \url{https://doi.org/10.1093/acprof:oso/9780199890712.003.0014}

\bibitem[\citeproctext]{ref-wellman2001metaanalysis}
Wellman, H. M., Cross, D., \& Watson, J. (2001). Meta-analysis of theory-of-mind development: {The} truth about false belief. \emph{Child Development}, \emph{72}, 655--684. \url{https://doi.org/10.1111/1467-8624.00304}

\bibitem[\citeproctext]{ref-zhang2021worse}
Zhang, Z., Yu, H., Long, M., \& Li, H. (2021). Worse {Theory} of {Mind} in {Only-Children Compared} to {Children With Siblings} and {Its Intervention}. \emph{Frontiers in Psychology}, \emph{12}, 5073. \url{https://doi.org/10.3389/fpsyg.2021.754168}

\end{CSLReferences}

\chapter{Appendix A --- Main Publications}\label{appendixA}

This dissertation includes four main publications that were either published (\hyperref[studyI]{Study I}) or under review (\hyperref[studyII]{Study II}, \hyperref[studyIII]{Study III}, \hyperref[studyIV]{Study IV}) at the time of the dissertation submission. The full texts of these publications are provided below. For the accepted manuscript, the published version is provided. For manuscripts under review, the submitted versions are provided which are published online as pre-prints.

~

\textbf{\hyperref[studyI]{Study I}:} Prein, J. C., Kalinke, S., Haun, D. B. M.*, \& Bohn, M.* (2023). TANGO: A reliable, open-source, browser-based task to assess individual differences in gaze understanding in 3 to 5-year-old children and adults. \emph{Behavior Research Methods, 56}(3), 2469--2485. \url{https://doi.org/10.3758/s13428-023-02159-5}

\textbf{\hyperref[studyII]{Study II}:} Prein, J. C., Maurits, L., Werwach, A., Haun, D. B. M.,* \& Bohn, M.* (2024). \emph{Variation in gaze understanding across the life span: A process-level perspective.} PsyArXiv. \url{https://doi.org/10.31234/osf.io/dy73a}

\textbf{\hyperref[studyIII]{Study III}:} Bohn, M.*, Prein, J. C.*, Ayikoru, A., Bednarski, F. M., Dzabatou, A., Frank, M. C., Henderson, A. M. E., Isabella, J., Kalbitz, J., Kanngiesser, P., Keşşafoğlu, D., Koymen, B., Manrique-Hernandez, M., Magazi, S., Mújica-Manrique, L., Ohlendorf, J., Olaoba, D., Pieters, W., Pope-Caldwell, S., \ldots{} Haun, D. (2024). \emph{A universal of human social cognition: Children from 17 communities process gaze in similar ways.} PsyArXiv. \url{https://doi.org/10.31234/osf.io/z3ahv}

\textbf{\hyperref[studyIV]{Study IV}:} Prein, J.C., Bednarski, F. M., Dzabatou, A., Frank, M. C., Henderson, A. M. E., Kalbitz, J., Kanngiesser, P., Keşşafoğlu, D., Köymen, B., Manrique-Hernandez, M. V., Magazi, S., Mújica-Manrique, L., Ohlendorf, J., Olaoba, D., Pieters, W. R., Pope-Caldwell, S., Sen, U., Slocombe, K., Sparks, R. Z., Stengelin, R., Sunderarajan, J., Sutherland, K., Tusiime, F., Vieira, W., Zhang, Z., Zong, Y., Haun, D. B. M.*, Bohn, M.*. (2024). \emph{Measuring variation in gaze following across communities, ages, and individuals -- a showcase of the TANGO--CC.} PsyArxiv.

\newpage

\section{Study I}\label{studyI}

\begin{minipage}{\textwidth}
\includepdf[pages={1}, scale=0.85, offset=0 -1cm, pagecommand={}]{../papers/studyI.pdf}
\end{minipage}

\includepdf[pages={2-}, scale=0.85, pagecommand={}]{../papers/studyI.pdf}

\newpage

\section{Study II}\label{studyII}

\begin{minipage}{\textwidth}
\includepdf[pages={1}, scale=0.85, offset=0 -1cm, pagecommand={}]{../papers/studyII.pdf}
\end{minipage}

\includepdf[pages={2-}, scale=0.85, pagecommand={}]{../papers/studyII.pdf}

\newpage

\section{Study III}\label{studyIII}

\begin{minipage}{\textwidth}
\includepdf[pages={1}, scale=0.85, offset=0 -1cm, pagecommand={}]{../papers/studyIII.pdf}
\end{minipage}

\includepdf[pages={2-}, scale=0.85, pagecommand={}]{../papers/studyIII.pdf}

\newpage

\section{Study IV}\label{studyIV}

\begin{minipage}{\textwidth}
\includepdf[pages={1}, scale=0.85, offset=0 -1cm, pagecommand={}]{../papers/studyIV.pdf}
\end{minipage}

\includepdf[pages={2-}, scale=0.85, pagecommand={}]{../papers/studyIV.pdf}

\chapter{Appendix B --- Further Publications}\label{appendixB}

The following Appendix B contains other publications that were written in the context of the dissertation but not included in the main text, with their respective abstracts.

\section{Action anticipation based on an agent's epistemic state in toddlers and adults}\label{action-anticipation-based-on-an-agents-epistemic-state-in-toddlers-and-adults}

\textbf{Citation:} Schuwerk, T., Kampis, D., Baillargeon, R., Biro, S., Bohn, M., Byers-Heinlein, K., Dörrenberg, S., Fisher, C., Franchin, L., Fulcher, T., Garbisch, I., Geraci, A., Grosse Wiesmann, C., Hamlin, K., Haun, D. B. M., Hepach, R., Hunnius, S., Hyde, D. C., Karman, P., \ldots, Prein, J., \ldots{} Rakoczy, H. (2021). \emph{Action anticipation based on an agent's epistemic state in toddlers and adults. Child Development} (In-Principle Acceptance of Registered Report Stage 1: Study Design). PsyArXiv. \url{https://doi.org/10.31234/osf.io/x4jbm}

\textbf{Abstract:} Do toddlers and adults engage in spontaneous Theory of Mind (ToM)? Evidence from anticipatory looking (AL) studies suggests that they do. But a growing body of failed replication studies raised questions about the paradigm's suitability. In this multi-lab collaboration, we test the robustness of spontaneous ToM measures. We examine whether 18- to 27-month-olds' and adults' anticipatory looks distinguish between two basic forms of an agent's epistemic states: knowledge and ignorance. In toddlers {[}ANTICIPATED n = 520 50\% FEMALE{]} and adults {[}ANTICIPATED n = 408, 50\% FEMALE{]} from diverse ethnic backgrounds, we found {[}SUPPORT/NO SUPPORT{]} for epistemic state-based action anticipation. Future research can probe whether this conclusion extends to more complex kinds of epistemic states, such as true and false beliefs.

\emph{Please note that this abstract was written for the Registered Report and does not entail results yet. Text in square brackets indicates placeholder text to be filled in after data collection.}

\newpage

\section{PREVIC: An adaptive parent report measure of expressive vocabulary in children between 3 and 8 years of age}\label{previc-an-adaptive-parent-report-measure-of-expressive-vocabulary-in-children-between-3-and-8-years-of-age}

\textbf{Citation:} Bohn, M., Prein, J. C., Engicht, J., Haun, D., Gagarina, N., \& Koch, T. (2023). \emph{PREVIC: An adaptive parent report measure of expressive vocabulary in children between 3 and 8 years of age.} PsyArXiv. \url{https://doi.org/10.31234/osf.io/hvncp}

\textbf{Abstract:} Parent report measures have proven to be a valuable research tool to study early language development. Caregivers are given a list of words and are asked which of them their child has already used. However, most available measures are not suited for children beyond infancy, come with substantial licensing costs or lack a clear psychometric foundation. Here we present the PREVIC (Parent Report of Expressive Vocabulary in Children), an open access, high quality vocabulary checklist for German-speaking children between three and eight years of age. The PREVIC was constructed leveraging the advantages of Item Response Theory: we designed a large initial item pool of 379 words and collected data from N = 1190 caregivers of children between three and eight years of age. Based on this data, we computed a range of fit indices for each item (word) and used an automated item selection algorithm to compile a final pool that contains items that a) vary in difficulty and b) fit the Rasch (one-parameter logistic) model. The resulting task is highly reliable and shows convergent validity. The IRT-based construction allowed us to design an adaptive version of the task, which substantially reduces the duration of the task while retaining measurement precision. The task -- including the adaptive version -- was implemented as a website and is freely accessible online (\url{https://ccp-odc.eva.mpg.de/previc-demo/}). The PREVIC fills an important gap in the toolkit of researchers interested in language development and provides an ideal starting point for the development of converging measures in other languages.

\newpage

\section{oREV: An item response theory-based open receptive vocabulary task for 3- to 8-year-old children}\label{orev-an-item-response-theory-based-open-receptive-vocabulary-task-for-3--to-8-year-old-children}

\textbf{Citation:} Bohn, M.*, Prein, J.*, Koch, T., Bee, R. M., Delikaya, B., Haun, D., \& Gagarina, N. (2024). oREV: An item response theory-based open receptive vocabulary task for 3- to 8-year-old children. \emph{Behavior Research Methods, 56}(3), 2595--2605. \url{https://doi.org/10.3758/s13428-023-02169-3}

\textbf{Abstract:} Individual differences in early language abilities are an important predictor of later life outcomes. High-quality, easy-access measures of language abilities are rare, especially in the preschool and primary school years. The present study describes the construction of a new receptive vocabulary task for children between 3 and 8 years of age. The task was implemented as a browser-based web application, allowing for both in-person and remote data collection via the internet. Based on data from N = 581 German-speaking children, we estimated the psychometric properties of each item in a larger initial item pool via item response modeling. We then applied an automated item selection procedure to select an optimal subset of items based on item difficulty and discrimination. The so-constructed task has 22 items and shows excellent psychometric properties with respect to reliability, stability, and convergent and discriminant validity. The construction, implementation, and item selection process described here makes it easy to extend the task or adapt it to different languages. All materials and code are freely accessible to interested researchers. The task can be used via the following website: \url{https://ccp-odc.eva.mpg.de/orev-demo}.

\newpage

\section{Validation of an open source, remote web-based eye-tracking method (WebGazer) for research in early childhood}\label{validation-of-an-open-source-remote-web-based-eye-tracking-method-webgazer-for-research-in-early-childhood}

\textbf{Citation:} Steffan, A., Zimmer, L., Arias-Trejo, N., Bohn, M., Dal Ben, R., Flores-Coronado, M. A., Franchin, L., Garbisch, I., Grosse Wiesmann, C., Hamlin, J. K., Havron, N., Hay, J. F., Hermansen, T. K., Jakobsen, K. V., Kalinke, S., Ko, E.-S., Kulke, L., Mayor, J., Meristo, M., \ldots, Prein, J., \ldots, Schuwerk, T. (2024). Validation of an open source, remote web-based eye-tracking method (WebGazer) for research in early childhood. \emph{Infancy, 29}(1), 31--55. \url{https://doi.org/10.1111/infa.12564}

\textbf{Abstract:} Measuring eye movements remotely via the participant's webcam promises to be an attractive methodological addition to in-person eye-tracking in the lab. However, there is a lack of systematic research comparing remote web-based eye-tracking with in-lab eye-tracking in young children. We report a multi-lab study that compared these two measures in an anticipatory looking task with toddlers using WebGazer.js and jsPsych. Results of our remotely tested sample of 18-27-month-old toddlers (N = 125) revealed that web-based eye-tracking successfully captured goal-based action predictions, although the proportion of the goal-directed anticipatory looking was lower compared to the in-lab sample (N = 70). As expected, attrition rate was substantially higher in the web-based (42\%) than the in-lab sample (10\%). Excluding trials based on visual inspection of the match of time-locked gaze coordinates and the participant's webcam video overlayed on the stimuli was an important preprocessing step to reduce noise in the data. We discuss the use of this remote web-based method in comparison with other current methodological innovations. Our study demonstrates that remote web-based eye-tracking can be a useful tool for testing toddlers, facilitating recruitment of larger and more diverse samples; a caveat to consider is the larger drop-out rate.

\chapter{Appendix C --- Social Cognition Survey}\label{appendixC}

In Autumn/Winter 2020, we conducted a short online expert survey on defining social cognition, as reported in the Introduction. Below you find the full survey, including the questions and the answers of the experts.

\section*{Documentation of the Expert Online Survey ``Defining Social Cognition''}\label{social-cognition-survey}
\addcontentsline{toc}{section}{Documentation of the Expert Online Survey ``Defining Social Cognition''}

\begin{minipage}{\textwidth}
\includepdf[pages={1}, scale=0.85, offset=0 -1cm, pagecommand={}]{../supplements/social_cognition_survey.pdf}
\end{minipage}

\includepdf[pages={2-}, scale=0.85, pagecommand={}]{../supplements/social_cognition_survey.pdf}

\newpage

\chapter{Selbstständigkeitserklärung}\label{selbststaendigkeit}

Julia Christin Prein\\
{[}Straße Hausnummer{]}\\
{[}PLZ Ort{]}\\
{[}Telefon{]}\\
{[}Email{]}\\

Hiermit erkläre ich, dass ich mich noch keiner Doktorprüfung unterzogen oder mich um Zulassung zu einer solchen beworben habe.

Ich versichere, dass die Dissertation \emph{Rethinking Variation in Social Cognition: Gaze Following across Individuals, Ages, and Communities} in der gegenwärtigen oder einer anderen Fassung noch keiner anderen Hochschule zur Begutachtung vorgelegen hat.

Ich versichere an Eides statt, dass ich die eingereichte Dissertation \emph{Rethinking Variation in Social Cognition: Gaze Following across Individuals, Ages, and Communities} selbstständig und ohne zulässige fremde Hilfe verfasst habe. Anderer als der von mir angegebenen Hilfsmittel und Schriften habe ich mich nicht bedient. Alle wörtlich oder sinngemäß anderen Schriften entnommenen Stellen habe ich kenntlich gemacht. Über die strafrechtlichen Folgen gemäß § 156 Strafgesetzbuch wurde ich in Kenntnis gesetzt.

~

Hamburg, {[}Datum{]}

~

{[}Unterschrift{]}

\end{document}
